\section{Introduction}
Periodic sensing and control traffic poses an energy--reliability tradeoff in
shared uplinks. A device can lower its transmit power and accumulate service
across several slots before its deadline. However, extending transmission also
increases the opportunities to interfere with other devices. With heterogeneous
periods and unaligned arrival phases, independently choosing a low power does
not ensure low network energy. Devices must also learn when to transmit and
which subchannel to use, including when to remain silent despite a pending
packet.

The energy benefit of stretching a transmission under deadline constraints is
established by classical packet-scheduling results~\cite{lazy2002}. In a shared
uplink, the practical question is how much of this benefit a distributed
controller can retain using delayed local feedback. High packet delivery alone
does not answer this question: frequent maximum-power transmissions may be
reliable but expensive. Conversely, low energy can result from leaving
difficult devices unserved. A useful evaluation must therefore combine
per-device reliability checks, an energy-aware local baseline, and an explicit
centralized reference.

Learning-based distributed spectrum access has been studied using recurrent
policies with local acknowledgments~\cite{naparstek2019}. Joint subband and
power decisions for delivering a payload within a time budget appear in
vehicular multi-agent reinforcement learning (MARL)~\cite{liang2019}.
For periodic mission-critical traffic, reinforcement-learning power control
has also been developed for reliability-constrained single-link
transmission~\cite{guo2023}. More recently, recurrent MARL has addressed spectrum
and power allocation with stochastic traffic in multicell
networks~\cite{zhang2025}. These works establish that neither joint resource
selection nor recurrent MARL is by itself a new contribution.

The particularly close XR-SBPC study~\cite{bagherinejad2025} combines periodic
traffic with offsets, accumulated Shannon service, and discrete subband--power
actions, including silence. It addresses reliability in a network of downlink
subnetworks and uses observations that include neighboring actions and current
link gains. Our study instead examines transmit-energy minimization under
per-device delivery requirements in a common-receiver uplink. Execution uses
only the device's own history and AP feedback, without access to other devices'
actions, queues, or phases. We evaluate this information model against
centralized references computed for the same packets and discrete power menu.

We instantiate a shared, subchannel-equivariant recurrent policy within
multi-agent proximal policy optimization (MAPPO)~\cite{yu2022}, and compare a
fixed violation penalty with local dual adaptation inspired by
state-augmented constrained reinforcement learning~\cite{calvo2021}.
The contribution is an evaluated access-control formulation and benchmark,
with three concrete components:
\begin{enumerate}
\item We specify asynchronous periodic delivery with accumulated service,
      joint subchannel--power--silence actions, and an explicit causal
      information boundary for decentralized execution.
\item We construct discrete-power zero-loss energy references: an exact
      interference-free solution when devices do not outnumber subchannels,
      and a separable lower reference plus a feasible cyclic orthogonal
      schedule at higher loads.
\item We quantify held-out energy and reliability with multiple training
      seeds, calibrated local rules, and paired uncertainty estimates.
      Additional diagnostics identify reduced transmission overlap as an
      observed correlate of the gains and expose the cost of insufficient
      training and residual per-device failures.
\end{enumerate}
Our conclusions concern static, frequency-flat channels and transmit energy.
They establish neither a general constraint guarantee for MAPPO nor superiority
over all available learning algorithms.
